\begin{table}[htbp]
\centering
\caption{Registro bruto da tabela \texttt{orgao.orgao} (SICOM, Caeté, 2023) antes da compactação do esquema}
\label{tab:registro_bruto}
\small
\begin{tabular}{ll}
\toprule
\textbf{Atributo (bruto)} & \textbf{Valor de exemplo} \\
\midrule
\texttt{seq\_orgao} & 304 \\
\texttt{num\_anoexercicio} & 2023 \\
\texttt{cod\_orgao} & 03 \\
\texttt{nom\_orgao} & SERVIÇO AUTÔNOMO DE ÁGUA E ESGOTO DE CAETÉ \\
\texttt{tipo\_orgao} & 03 - AUTARQUIA (EXCETO RPPS) \\
\texttt{cod\_municipio}* & 3110004 \\
\texttt{nom\_municipio}* & CAETÉ \\
\texttt{cod\_uf} & 31 \\
\texttt{sgl\_uf} & MG \\
\texttt{nom\_uf} & Minas Gerais \\
\texttt{dsc\_regiaoplanejamento} & CENTRAL \\
\bottomrule
\end{tabular}
\vspace{2pt}
\begin{minipage}{0.95\linewidth}
\footnotesize *Atributos removidos na compactação do esquema (Seção 3.3): a base é exclusiva do município de Caeté, portanto o código e o nome do município são redundantes em toda tabela que os contém --- assim como \texttt{cod\_uf}, \texttt{sgl\_uf} e \texttt{nom\_uf} nesta mesma tabela, que também não variam entre registros mas não são removidos pela regra de compactação por não fazerem parte da lista de atributos de controle identificada no pipeline ETL. A mesma poda remove ainda, em outras tabelas do esquema não ilustradas aqui, os atributos de controle de arquivo/lote \texttt{num\_versao\_arq} (tabelas de frota) e \texttt{num\_lote} (tabelas de licitação).
\end{minipage}
\end{table}
